\newpage
\phantomsection
\label{prf:lebesgue-criterion-riemann-integrability}
\LRAProofFor{thm:lebesgue-criterion-riemann-integrability}
\begin{remark*}[Return]\hyperref[thm:lebesgue-criterion-riemann-integrability]{Return to Theorem}\end{remark*}
\begin{theorem*}[Lebesgue Criterion for Riemann Integrability]
A bounded function \(f:[a,b]\to\mathbb R\) is Riemann-integrable if and only
if its discontinuity set has measure zero.
\end{theorem*}
\begin{proof}[Professional Standard Proof]
\LRAProofBodyStart
Let
\[
  \Omega=\sup_{[a,b]}f-\inf_{[a,b]}f<\infty.
\]

First suppose the discontinuity set \(D\) has measure zero. Fix
\(\varepsilon>0\). Choose \(\eta>0\) with \(\eta(b-a)<\varepsilon/2\). Let
\[
  D_\eta=\{x:\omega_f(x)\ge \eta\}.
\]
Then \(D_\eta\subseteq D\), so \(D_\eta\) has measure zero. Cover \(D_\eta\)
by open intervals whose total length is less than \(\varepsilon/(2\Omega)\);
by compactness of \(D_\eta\), take a finite subcover and call its union \(G\).
On \(K=[a,b]\setminus G\), every point has oscillation \(<\eta\). Compactness
of \(K\) gives a single scale \(\delta>0\) such that every interval of length
less than \(\delta\) meeting \(K\) has oscillation \(<\eta\).

Choose a partition \(P\) with mesh \(<\delta\) that includes the endpoints of
the intervals making up \(G\). Slabs contained in \(G\) have total width less
than \(\varepsilon/(2\Omega)\), so their contribution to \(U(f,P)-L(f,P)\) is
less than \(\varepsilon/2\). Slabs meeting \(K\) have oscillation \(<\eta\), so
their total contribution is less than \(\eta(b-a)<\varepsilon/2\). Hence
\[
  U(f,P)-L(f,P)<\varepsilon.
\]
By the Darboux criterion, \(f\) is Riemann-integrable.

Conversely, suppose \(D\) does not have measure zero. Since
\[
  D=\bigcup_{k=1}^{\infty}D_{1/k},
\]
some \(D_{\eta_0}\) with \(\eta_0=1/k_0\) does not have measure zero. Thus
there is \(c>0\) such that every interval cover of \(D_{\eta_0}\) has total
length at least \(c\). For any partition \(P\), the slabs whose interiors meet
\(D_{\eta_0}\), together with the finitely many partition points, cover
\(D_{\eta_0}\); ignoring finitely many points does not affect the lower bound
on total interval length. Hence those slabs have total width at least \(c\).
Each such slab has oscillation at least \(\eta_0\), so
\[
  U(f,P)-L(f,P)
  =\sum_i \omega_f([x_{i-1},x_i])\Delta x_i
  \ge \eta_0 c>0.
\]
The Darboux gap cannot be made arbitrarily small, so \(f\) is not integrable.
\end{proof}
\begin{proof}[Detailed Learning Proof]
\LRAProofBodyStart
The proof says that all Darboux error lives over points where the function
oscillates. If those bad points fit inside intervals of arbitrarily small total
length, their worst possible contribution is small, and the good complement
has uniformly small oscillation. If the bad set has positive unavoidable
length at some oscillation level, every partition pays a fixed positive
Darboux gap.
\end{proof}
\begin{remark*}[Proof structure]
Use the sets \(D_\eta=\{x:\omega_f(x)\ge\eta\}\) to separate large oscillation
from small oscillation, then apply the Darboux criterion in both directions.
\end{remark*}
\begin{dependencies}
\begin{itemize}
  \item \hyperref[def:measure-zero]{Measure Zero}.
  \item \hyperref[def:point-oscillation-integration]{Point Oscillation}.
  \item \hyperref[thm:darboux-criterion]{Darboux Criterion}.
\end{itemize}
\end{dependencies}
\clearpage
